\chapter{Reproducibility Records}
\label{app:reproducibility}

\section{Source, environment and build}

The project repository is identified by \citep{mfjarepository}. The reference
environment is Ubuntu~24.04, ROS~2 Jazzy, Gazebo Harmonic and Python~3.12. The
revision carrying this report stores the implementation, active and recovery
configurations, tests, lightweight evidence and submission PDF together. A
clean checkout of that revision therefore reconstructs the source handover.
The model checkpoint, sealed dataset and raw ROS bags are external
objects identified by the fingerprints below rather than files stored in Git.
Here, a checkpoint is a saved network-weight state, a sealed dataset is a
frozen hash-identified dataset (not an encrypted or still-unseen one), and a ROS
bag is a timestamped recording of ROS messages. An external object is retained
outside Git because of size, while its digest fixes the object used by the
reported run.

From the workspace root, the source revision can be recorded and the installed
overlay rebuilt with the commands below. In ROS/colcon terminology, an overlay
is this workspace's build and install output layered over the base ROS
installation.

\begin{lstlisting}[language=bash]
source /opt/ros/jazzy/setup.bash
git -C src/mfja_3rd_floor_gz rev-parse HEAD
git -C src/mfja_3rd_floor_gz status --short
colcon build --symlink-install --base-paths src/mfja_3rd_floor_gz
source install/setup.bash
\end{lstlisting}

From \path{report/}, \code{make check} builds the A4 PDF, rejects unresolved
references or citations and checks for overfull boxes. The retained build log
and delivery-PDF fingerprint are stored under \path{report/evidence/}; keeping
the PDF digest outside the PDF avoids a self-referential identity. Operating
procedures and launch arguments are documented in \path{runbook.html},
\path{docs/INSTALLATION.md}, \path{docs/QUICK_START_AND_FEATURE_GUIDE.md} and
\path{docs/ROOM315_LANGUAGE_TO_MOTION_RUNTIME.md}.

\section{Primary evidence identity}

The repository-relative evidence package is
\path{report/evidence/room315_visual_v4_submission_2026-08-11}. Its
\path{evidence_manifest.json} has SHA-256
\path{840787b617e0f671628dfc7d8122ff559d76664506ff8f94ac31d043737da446},
and the package-wide \path{SHA256SUMS} file has SHA-256
\path{dec74565a8ccfba57a32d83dfdd03236daeafeff226d5243594e328fa4adc5ab}.
Together they cryptographically bind the lightweight copies and external-object
fingerprints, and map each reported result to its source record. The complete
per-file hashes and JSON Pointers remain
in that package and in \path{report/evidence/ACTIVE_EVIDENCE.md}; they are not
repeated in this appendix.

\emph{Lightweight} means that the package retains manifests, controls, results
and checksums but not the large checkpoint, images or raw bags. Binding means
that any content change produces a different SHA-256 and fails verification; it
does not make a file physically impossible to edit. A JSON Pointer is the
recorded path to one field inside a JSON document.

The selected epoch-11 visual checkpoint has SHA-256
\path{869d64049b0092c37d21a4c8b910dc6b91954527e0e49c5694fa82dce570f40d}.
The sealed 1,040-scene Final Test has dataset fingerprint
\path{226f4b5889f7d8b66bccc87d3e8dd494f710c8f4b228fab4851d6da7448e2eef}.
These fingerprints identify the two principal external objects used by the
visual evaluation; each is the SHA-256 identifier recorded by its governing
protocol. Some copied machine-generated records preserve absolute
archival paths from the evaluation host. They are retained for provenance but
are not submission locations; the authored evidence index uses relative paths.

\section{Result-to-record index}

\Cref{tab:evidence-roots} gives the shortest path from each reported result to
its inspectable record. Paths beginning with \path{model/}, \path{final_test/}
or \path{runtime/} are relative to the primary evidence package above.

\begin{longtable}{@{}
>{\raggedright\arraybackslash}p{0.22\textwidth}
>{\raggedright\arraybackslash}p{0.67\textwidth}@{}}
\caption{Records supporting the reported results.}
\label{tab:evidence-roots}\\
\toprule
\textbf{Result group} & \textbf{Records and scope}\\
\midrule
\endfirsthead
\toprule
\textbf{Result group} & \textbf{Records and scope}\\
\midrule
\endhead
Visual model and Final Test & \path{model/} contains the effective training
configuration, input fingerprints, selected-run report and hard-case Validation records;
\path{canary/} is the labelled development regression check.
\path{final_test/final/controls/} contains the sealed controls and one-shot
evaluation contract, while \path{final_test/final/results/} contains the
immutable completion ledger and measurements. The controls freeze what may run;
the ledger records that the single permitted attempt was consumed and
completed.\\
\tablerowrule
Runtime qualification & \path{runtime/shadow/}, \path{runtime/campaign/} and
\path{runtime/safety/} contain the observation-only comparison, scene campaign
and visual-input fault checks. \path{runtime/closed_loop/positive/},
\path{runtime/closed_loop/fault/} and
\path{runtime/closed_loop/active_runtime/} contain the positive campaign, fault
campaign, manually approved Gazebo runtime and its selected-case smoke. Exact
per-record digests are available in
\path{report/evidence/ACTIVE_EVIDENCE.md}.\\
\tablerowrule
Language interface &
\path{report/evidence/room315_language_semantic_evaluation_v3} contains the
pinned ten-case corpus, configuration, raw model envelopes, fusion traces,
environment and checksums. The directory suffix is an evidence revision and
does not select the visual runtime.\\
\tablerowrule
Software and full-floor checks & The recorded component commands and outputs
are the three \path{pytest_current_source_*_2026-08-07.log} files under
\path{report/evidence/}. The isolated all-robot launch/interface record is
\path{report/evidence/multi_robot_cold_start_smoke_2026-08-07}.\\
\tablerowrule
Experimental camera figure & The exact B03 frames, extraction program,
provenance and checksums for \cref{fig:b03-camera-evidence} are under
\path{report/evidence/room315_v4_closed_loop_campaign_figures_2026-08-12}.\\
\tablerowrule
Runtime and recovery configuration & The active files are
\path{mfja_robot_control_config/config/room_315_vla/visual_state_runtime.yaml}
and \path{task_execution_runtime.yaml}. Explicit recovery profiles are
\path{visual_state_runtime_v3_rollback.yaml} and
\path{task_execution_runtime_v3_rollback.yaml}. The task configuration defaults
to \code{execution\_enabled=false}. These recovery profiles are explicit manual
fallback configurations, not automatically selected alternatives.\\
\bottomrule
\end{longtable}
